\documentclass[11pt]{article}
\usepackage[margin=1in]{geometry}
\usepackage{amsmath,amssymb,amsthm,booktabs,hyperref}
\newtheorem{theorem}{Theorem}
\newtheorem{proposition}[theorem]{Proposition}
\theoremstyle{remark}\newtheorem*{remark}{Remark}
\newcommand{\F}{\mathcal{F}}
\title{Strong unextendibility of small Pauli MUB sets:\\ exact bounds for two, three and four qubits}
\author{Draft (W33 repository, Passes 12636--12648)}
\date{November 2026}
\begin{document}
\maketitle
\begin{abstract}
A set of mutually unbiased stabiliser bases (Pauli MUBs) is \emph{weakly unextendible} if no further Pauli class is unbiased to all of
them, and \emph{strongly unextendible} if no vector at all is. For a set $\mathcal{S}$ in dimension $d$ we study
$\F(v)=\sum_{e}\bigl(|\langle e|v\rangle|^2-|v|^2/d\bigr)^2$, the sum over all basis vectors $e$ of the set; strong unextendibility
is $\min_{|v|=1}\F>0$. We prove exact, quantitative lower bounds. For two qubits, the minimum is exactly $1/4$. For three qubits, all
24\,192 five-element sets form one Clifford orbit and satisfy $\F\ge(1/8-10^{-6})|v|^4$, while stabiliser states attain $1/8$. For
four qubits, the smallest weakly unextendible sets have exactly five elements (an exhaustive search, matching Grassl's size
$q^{n/2}+1$). There are $105\,283\,584$ of them in two types. One type satisfies $\F\ge|v|^4/80$ exactly, so five bases are strongly
unextendible in $d=16$, and its minimum ($\approx0.0524$) lies below every stabiliser value. The proofs combine a twirl identity, the
graph-theoretic uncertainty relation of de~Gois, Hansenne and G\"uhne, Pauli-symmetry reduction of sum-of-squares certificates, facial
reduction, and verification in exact rational arithmetic.
\end{abstract}

\section{Setting}
Let $n$ qubits carry the Pauli group with symplectic space $W(2n-1,2)$. Its points are Pauli operators up to phase, and its
Lagrangians are maximal commuting classes. Disjoint Lagrangians give mutually unbiased stabiliser bases, and a weakly unextendible
set is a maximal partial spread~\cite{mbgw,thas}. Mandayam, Bandyopadhyay, Grassl and Wootters~\cite{mbgw} constructed such sets with
$d/2+1$ elements for $d=4,8$. They proved strong unextendibility for specific examples with a Gr\"obner basis computation, and
noted that a general proof was open. Kantor~\cite{kantor} surveys maximal symplectic partial spreads and reports Grassl's sets of
size $q^{n/2}+1$.

\paragraph{Twirl identity.} For a stabiliser basis $B_L$, $\sum_{e\in B_L}|e\rangle\langle e|^{\otimes2}=d^{-1}\sum_{P\in L}P^{\otimes2}$.
Hence, with $C$ the covered Paulis and $U$ the uncovered ones,
\begin{equation}
\F(v)=\frac1d\sum_{P\in C}\langle P\rangle^2=\frac1d\Bigl(d-1-\sum_{P\in U}\langle P\rangle^2\Bigr)\qquad(|v|=1),
\end{equation}
since $\sum_{P\ne I}\langle P\rangle^2=d-1$ for pure states. A stabiliser state of a Lagrangian meeting $C$ in $m$ points gives
$\F=m/d$. Maximality forces $m\ge1$, so $\min\F\le m^*/d$, where $m^*$ is the least such $m$.

\paragraph{Lovász bound.} De~Gois, Hansenne and G\"uhne~\cite{dhg} prove $\sum_i\langle A_i\rangle^2\le\vartheta(G)$ for dichotomic
observables, where $G$ is their anticommutation graph and $\vartheta$ the Lovász number~\cite{lovasz}. Applied to $U$,
$\min\F\ge(d-1-\vartheta)/d$. The bound is tight when $\vartheta=\alpha(G)=d-1-m^*$.

\section{Results}
\begin{theorem}[two qubits]\label{t:2}
For every three-element maximal partial spread of $W(3,2)$, $\min\F=1/4$ exactly. Here $\vartheta=\alpha=2$, and equivalently the six
uncovered Paulis split into two pairwise-anticommuting triples.
\end{theorem}

\begin{theorem}[three qubits]\label{t:3}
The maximal partial spreads of $W(5,2)$ have sizes 5 (24\,192 sets, one $\mathrm{Sp}(6,2)$-orbit) and 9 (complete sets). For every
five-element set and every $v\in\mathbb{C}^8$, $\F(v)\ge(1/8-10^{-6})|v|^4$, and stabiliser states attain $1/8$. A weaker but
structural bound is $\F\ge\frac{3-\sqrt5}{8}|v|^4$, from $\vartheta\le4+\sqrt5$, certified over $\mathbb{Q}(\sqrt5)$ in the rank-21
coherent closure of the uncovered graph.
\end{theorem}
\noindent\emph{Method.} $\F(v)|v|^2-\delta|v|^6$ is written as an SOS in the degree-3 monomials and in the monomials
$v_iv_j\bar v_k$. The Pauli group acts on both monomial spaces with centre $-1$, so each is a multiple of the unique 8-dimensional
irrep, and invariant Gram matrices reduce to blocks of size 15 and 36. All coefficients are Gaussian rationals. At
$\delta=1/8-\varepsilon$, a strictly interior certificate is rounded (denominator 16). Its pivot variables are restored exactly, every
coefficient equation is verified in integer arithmetic, and positivity is checked by exact $LDL^{\mathsf T}$, for
$\varepsilon=10^{-3},\dots,10^{-6}$. At $\delta=1/8$ the 160 stabiliser minimisers force the 15-block to vanish, and facial reduction
leaves a 16-dimensional face with zero interior margin. This suggests that the degree-6 bound approaches $1/8$ without attaining it.

\begin{theorem}[four qubits]\label{t:4}
(i) $W(7,2)$ has no maximal partial spread of size 3 or 4. The smallest size is 5, and there are $105\,283\,584$ such sets.
(ii) They fall into two invariant types. For a set of the first type (lines listed in the repository), $\vartheta\le74/5$, so
$\F\ge|v|^4/80$ exactly, while numerically $\min\F=0.0523626<1/16$. Sets of the second type are \emph{lifts} of the three-qubit five-sets. For each
of them, 120 uncovered points $p$ project the set onto a maximal five-set of $p^\perp/\langle p\rangle\cong W(5,2)$, whereas the first
type has no such point. They inherit a mirror partner (a second partition of the same covered set) and have numerical minimum $1/16$.
(iii) For several inequivalent nine-element sets, $\min\F=m^*/16$ exactly ($1/16$ or $1/8$), via a facially reduced Lovász dual.
\end{theorem}
\noindent\emph{Method for (i).} By Witt's theorem, $\mathrm{Sp}(8,2)$ is transitive on ordered pairs of complementary Lagrangians, so
every candidate can be taken to contain a fixed pair. The remaining one or two Lagrangians are enumerated exhaustively. As a control,
the same code finds no sizes 3 and 4 for three qubits, and finds 56 five-sets through a fixed pair, matching
$24\,192\cdot20/(135\cdot64)=56$. \emph{Method for (iii).} The maximum independent sets of the uncovered graph are the Lagrangians
meeting $C$ in $m^*$ points, minus those points. Complementary slackness with them gives linear equations on the dual. Imposing these
equations makes rounding succeed, and the certificate is checked by exact $LDL^{\mathsf T}$.

\begin{remark}[two qutrits]
For the 432 five-element maximal partial spreads of $W(3,3)$, $\F\ge|v|^4/9$ exactly, via an $S_6$/Jucys--Murphy identity. The true
minimum is numerically $(207-19\sqrt{57})/432\approx0.14711$, which is an identification, not a proof. Mirror partners there lie in
opposite orbits and are exchanged only by antiunitaries; a third-moment witness of size $16/27$ separates them. In characteristic 2
the corresponding partners are Clifford-equivalent.
\end{remark}

\section{Summary and open questions}
\begin{center}
\begin{tabular}{llll}\toprule
$d$ & smallest size & $\min\F$ & status\\\midrule
4 & 3 & $1/4$ & exact\\
8 & 5 & $\in[1/8-10^{-6},\,1/8]$ & exact interval\\
9 & 5 & $\ge1/9$; $\approx0.14711$ & exact bound; numerical value\\
16 & 5 & $\ge1/80$; $\approx0.05236$ (one type) & exact bound; numerical value\\\bottomrule
\end{tabular}
\end{center}
Open: an exact certificate at $1/8$ (degree 8, or a structural proof); closed forms for the non-stabiliser optima; the strong
unextendibility constant for all five-element sets in $d=16$ and beyond. \emph{Not claimed}: novelty of the weak-unextendibility
census at $d=16$, which may already be in Grassl's tables~\cite{kantor}. All certificates are in the repository (\texttt{data/*.json})
with the scripts that verify them.

\begin{thebibliography}{9}
\bibitem{mbgw} P.~Mandayam, S.~Bandyopadhyay, M.~Grassl, W.~K.~Wootters, Unextendible mutually unbiased bases from Pauli classes,
Quantum Inf.\ Comput.\ 14 (2014); arXiv:1302.3709.
\bibitem{thas} K.~Thas, Unextendible mutually unbiased bases (after Mandayam, Bandyopadhyay, Grassl and Wootters), arXiv:1407.2778.
\bibitem{kantor} W.~M.~Kantor, On maximal symplectic partial spreads, Adv.\ Geom.\ (2017).
\bibitem{dhg} C.~de~Gois, K.~Hansenne, O.~G\"uhne, Uncertainty relations from graph theory, Phys.\ Rev.\ A 107, 062211 (2023).
\bibitem{lovasz} L.~Lov\'asz, On the Shannon capacity of a graph, IEEE Trans.\ Inf.\ Theory 25 (1979).
\end{thebibliography}
\end{document}
